\section{Optional Hardening Tier}
\label{sec:hardening}

\subsection{Why a separate tier}

The safe tier deliberately forbids \code{unsafe}. Debugger detection, timing
checks, debug-register inspection, and code-page hashing require OS interfaces,
privileged instructions, or raw memory access that safe Rust cannot express. The
project keeps the default artifact safe and places these controls in an opt-in
\code{sirus-hardened} crate, enabled only with \verb|--hardened|. The hardened
crate is never linked into the default runtime, and the safe tier's
\code{forbid(unsafe_code)} policy remains enforced by the build.

This separation is a security property. A user who does not need the hardened
tier does not inherit its unsafe code, its platform-specific failure modes, or
its compatibility risk.

\subsection{Linux tracer detection}

On Linux, the hardened entry reads \code{/proc/self/status} through a raw
syscall and checks the \code{TracerPid} field. A non-zero tracer PID causes the
entry to fail closed with value zero. The check is intentionally narrow: it
detects a class of ptrace-based debuggers and does not claim to detect kernel
modules, hypervisors, or all dynamic instrumentation.

The raw syscall avoids a libc dependency and keeps the generated runtime
freestanding. The implementation is unsafe by construction; it is confined to
the hardened crate and reviewed as a boundary component.

\subsection{Windows debugger checks}

On Windows, the hardened entry inspects the process environment block for
\code{BeingDebugged} and \code{NtGlobalFlag}. A detected debugger causes a
fail-closed return. This is a user-mode check; it does not claim to detect a
kernel debugger or a hypervisor.

The Windows debug-register path is available through the \code{harden-check}
command. It uses the thread context API to report the debug registers on a native
Windows host. Linux user mode cannot read debug registers without a ptrace
relationship, so the command reports \code{unsupported} there rather than
pretending to provide a control.

\subsection{RDTSC and LFENCE timing budget}

The hardened entry measures a bounded interval with \code{rdtsc}, serialized with
\code{lfence}, and fails if the interval exceeds the configured budget. This
raises the cost of single-stepping and some emulation-based tracing. It is not a
proof against a hypervisor that controls the timestamp counter, and it can
produce false positives on a heavily loaded or migrated process.

The check is opt-in because timing checks interact with scheduling, virtualization,
and power management. The safe tier does not include them.

\subsection{Linked text checksum}

Linux and Windows hardened builds compute a checksum over the linked VM text at
runtime and compare it with the build-time value. A modified VM text region
causes the entry to fail closed. The checksum is implemented in the hardened
crate with raw memory reads; it is not present in the safe tier.

The checksum protects the VM text, not the entire host image. A patcher that
changes the host entry, the payload, or the expected checksum is outside its
scope. The tamper test modifies VM text and verifies that the protected binary
rejects the modification.

\subsection{Failure behavior}

A failed hardened check returns zero. As with all fail-closed paths, the intended
deployment consumes the protected result as data, so a zero result causes a
downstream failure rather than a silent success. The generated entry does not
abort, print, or allocate.

\subsection{Boundaries}

The hardened tier does not provide kernel-level anti-debug, hypervisor detection,
DBI detection, trusted signing, or a formal Spectre mitigation. Those controls
require privileges, kernel components, or hardware features outside the project's
safe-Rust and local-evidence constraints. The paper reports them as
non-capabilities, not as hidden features.
